% !TEX root = 07-orchestration.tex
% title: 协议编排

笔记 00 $\sim$ 06 备齐了全部零件：
Endemic OT 产出 Base OT 种子，PPRF 和 SoftSpokenOT 把种子扩展成海量随机 OT，
随机 VOLE 再把随机 OT 替换成 MtA 份额。
本篇把零件装配成完整的门限 ECDSA。

【协议参数】

\begin{itemize}
\item
  参与方集合 $\mathbb{K}$，元素 $i\in\mathbb{Z}_n^*$，一般为从 1 开始的连续整数。
\item
  参与方编号 $i \ge 1$ 同时也是 Lagrange 插值点，视作 $\mathbb{Z}_n^*$ 元素。
\item
  阈值 $t$，满足签名参与方数量 $\left| \mathbb{S} \right| \ge t$。
\item
  $\mathtt{sid}$ 是协议外部输入，全员事先约定。
\item
  子协议会话号从 $\mathtt{sid}$ 和参与方编号派生而来。
  例如 Base OT 实例 $(i,j)$ 使用会话号
\end{itemize}
\begin{equation}
\label{eq:orch:subsid-example}\tag{subsid-example}
\mathrm{Hash}(\mathtt{sid}, i, j, \texttt{"base\_ot"}). 
\end{equation}

【变量命名约定】$i$ 是本方编号，$j\ne i$ 是其他参与方编号。后文默认遵循这个约定，除非特别说明。
遇到动词“收发”、“交换”时，以明确写进宾语的变量作为本方发送的消息。
对宾语中的变量名进行变换「把下标 $i$ 改名为 $j$、$j$（如果有）改名为 $i$」，
隐式声明这些变量为本方接收的消息。

【检查级别与失败处理备查】
\label{sec:orch:failure-levels}

本文用 K、S、P 区分失败的作用范围，并在相关段落开头标注【检查-K】、【检查-S】或【检查-P】。
这些字母是本文的编排记号，不是 DKLS23 原文的错误码，也不表示失败已经证明某方作弊。
各类失败都立即停止受影响的执行，向调用方报告失败，销毁该执行的临时秘密状态，且不输出其结果。
通过认证信道向本次参与方发送绑定会话号的失败通知；通知中不包含秘密份额、OT 密钥或种子。
通知是终止信号，不是可转交的作弊证据；收到通知的一方不得仅凭通知发送者的身份将其停用。

\begin{description}
\item[【检查-K】Keygen 失败。]
  适用于初始化中的承诺、DLog 证明、Feldman 份额、Endemic OT 与 PPRF 检验。
  中止本次 Keygen，销毁本次生成的秘密材料，不输出或使用未完成的 Keyshare。
  拒绝揭示、揭示长度不符，以及收到本次 Keygen 的失败通知，也按此处理。
  若重新生成密钥，使用新的 Keygen 会话号和全新的初始化随机数。
\item[【检查-S】单次签名失败。]
  中止本次 PreSign 或 Sign，废弃对应预签名及临时材料；保留既有 Keyshare。
  适用于总点为无穷远点、$r$ 或 $s$ 无效、签名分母为零、最终验签失败，
  以及收到本次签名的失败通知；仅凭这些结果不新增停用对端。
  若重试，须新建 PreSign 会话并重新采样，不能复用失败的预签名。
\item[【检查-P】停用对端并终止关联会话。]
  先执行 S 类处理，再将涉及的对端集合 $B$ 记入本方针对当前 Keyshare 的持久停用记录。
  若本地失败可归于与 $j$ 的具体交互，则 $B:=B\cup \{j\}$；这表示停止交互，不等于公开归责。
  对本方所有涉及 $B$ 中任一方的并发签名会话，停止继续计算并在相应协议位置发送失败通知。
  同时废弃本方所有涉及这些对端的待用预签名，拒绝后续与这些对端共同签名。
  不涉及这些对端的会话可以继续；更换 PreSign 会话号不能清除停用记录。
\end{description}

P 类处理以 DKLS23 Protocol 3.6 第 8 步为依据，涵盖 RVOLE 失败与逐对 $\Gamma$ 等式失败。
第 9 步规定收到失败通知时终止本次签名，第 10 步规定最终验签失败时报告本次失败，均归入 S 类。
原文第 8 步还要求聚合公钥一致，但这个聚合等式本身不能定位某个对端。
本文对无法定位的聚合公钥或 transcript 不一致，保守取 $B=\mathbb S\setminus\{i\}$；
这会连带停用无辜对端，是本文明确采用的可用性取舍，不能据此声称已识别作弊者。
具体实现中的 SoftSpoken 验证失败与 nonce 承诺揭示失败，也采用 P 类处理。
这些实现层面的停用范围是本文补充的保守规则，不应表述为原文逐项给出的归责机制。
原文参见\href{https://eprint.iacr.org/2023/765.pdf}{第 3 节，Protocol 3.6}。

等待消息本身不等于检验失败：异步网络中的延迟不能证明对方拒绝发送。
若调用方因超时取消，Keygen 按 K 类、签名按 S 类收尾，不据此新增停用对端。
S 类重试必须同时服从已有的 P 类停用记录；本文不提供清除记录、重新接纳对端的恢复流程。

\section{分布式密钥生成（DKG、Keygen）}

【与原文 Keygen 的关系】本节记录的是采用 Feldman 承诺与 DLog 知识证明的 Keygen 流程，
并非 DKLS23 Protocol 7.1 的逐步转写。
原文的 relaxed Keygen 保留公钥与份额的一致性及门限私钥保护，
放宽的是安全证明对提取恶意方私钥份额的要求，因此其协议不使用知识证明。
DKLS23 在签名阶段重新随机化份额，并通过 RVOLE 提取证明所需的输入，因而可以采用这种 Keygen。
本节流程来自笔记整理时参考的实现与材料；具体来源尚未逐项对应，不能直接套用 Protocol 7.1 的证明。
相关说明见 \href{https://eprint.iacr.org/2023/765.pdf}{DKLS23 \S3.1、\S7 与 Protocol 7.1}。

\subsection{规格}

【输入】公开输入：参与方集合 $\mathbb{K}$，签名所需最少参与方数量（门限） $t$，
本方编号 $i\in\mathbb{K}$，会话编号 $\mathtt{sid}$。

【输出】每方 $i$ 获得 Keyshare，内含：

\begin{itemize}
\item
  公钥 $Y=xG$，全员一致；
\item
  编号 $i$，同时也是多项式的求值点；
\item
  私钥分片 $x_i := P(i)$，即全局秘密多项式 $P$ 在自己编号处的求值，
  满足 $x=\sum_{j\in\mathbb{S}}\lambda(j,\mathbb{S})\cdot x_j$，插值系数见式 \eqref{eq:orch:lagrange-coef}；
\item
  针对每个 $j\ne i$ 以及树编号 $h\in[0,\kappa/k)$，
  保存本方作为 PPRF Sender 的 $\mathcal{A}_{h,i,j}$，
  以及作为 PPRF Receiver 的 $\bigl(y_{h,i,j}, \mathcal{A}^*_{h,i,j}\bigr)$，
  供 Sign 阶段的 SoftSpoken 取用；
\item
  针对每个 $j\in\mathbb{K}\setminus\{i\}$，
  保存零和掩码分片 $\epsilon_{i,j}$，供协议 PreSign 派生零和掩码。
  两方保存同一个值，也就是说 $\epsilon_{i,j}=\epsilon_{j,i}$。
\end{itemize}

【安全承诺】

\begin{itemize}
\item
  不足 $t$ 方的联盟对私钥 $x$ 一无所知；
\item
  【检查-K】hash 承诺、DLog 证明、Feldman 与 PPRF 验证失败时，执行 K 类处理。
\end{itemize}

\subsection{实施}

\subsubsection{Feldman 承诺与零和掩码种子承诺}
\label{sec:orchestration:round1}

本方生成一个 $t-1$ 次随机多项式 $P_i \in \mathbb{Z}_n[X]$，
如公式 \eqref{eq:orch:poly-share} 所示。
各方独立均匀采样全部系数，包括常数项。
采样完成后，以各方常数项之和定义私钥 $x$，如公式 \eqref{eq:orch:poly-c0}；$x$ 不是预先给定的约束。
\begin{align}
\tag{coeff}
\phantom{{}={}} c_{i,*} & \stackrel{\$}{\leftarrow} \mathbb{Z}_n, \\
\label{eq:orch:poly-share}\tag{PiX}
P_i(X) &:= c_{i,0} + c_{i,1}X + \cdots + c_{i,t-1}X^{t-1}, \\
\label{eq:orch:poly-c0}\tag{c0}
x &:= \sum_{i\in\mathbb{K}} c_{i,0} \pmod n.
\end{align}

记 Feldman 向量 $\mathbf{F}_i$ 如公式 \eqref{eq:orch:feld-share}。
本方对 $\mathbf{F}_i$ 进行哈希承诺如公式 \eqref{eq:orch:feld-com}，
式中 $\varepsilon_{1,i} \stackrel{\$}{\leftarrow} \left\{0,1\right\}^{256}$ 是盲化项。
\begin{align}
\label{eq:orch:feld-share}\tag{Fi}
\mathbf{F}_i
&:= (c_{i,0} \cdot G, \cdots, c_{i,t-1}\cdot G). \\
\label{eq:orch:feld-com}\tag{Com1}
\mathrm{Com}_{1,i}
&:= \mathrm{Hash}\left(
  \mathrm{sid}, i, \mathbf{F}_i, \varepsilon_{1,i}
\right).
\end{align}

【检查-K】本方还要为每个 $j\in\mathbb{K}\setminus\{i\}$ 独立采样种子贡献 $\epsilon^\circ_{i,j}$，
以及用于隐藏承诺内容的随机串 $\nu_{i,j}$，两者均为 256 比特。
这里的 $\epsilon^\circ_{i,j}$ 属于本方，$\epsilon^\circ_{j,i}$ 属于对方，二者无需相等。
它们将在揭示并验证后合成为同一个共享种子 $\epsilon_{i,j}$。
\begin{align}
  \label{eq:orch:seed-contribution}\tag{seed-contribution}
  \epsilon^\circ_{i,j},\nu_{i,j}
  &\stackrel{\$}{\leftarrow}\{0,1\}^{256}, \\
  \label{eq:orch:seed-com}\tag{seed-com}
  \mathrm{Com}_{\epsilon,i,j}
  &:=\mathrm{Hash}\left(
    \mathrm{sid},i,j,\texttt{"zero\_seed"},\epsilon^\circ_{i,j},\nu_{i,j}
  \right).
\end{align}

【交换】 $\mathrm{Com}_{1,i}, \mathrm{Com}_{\epsilon,i,j}$。

\subsubsection{启动 Endemic OT，揭示 Feldman 与种子承诺}

本方为 $P_i$ 的每个系数制作 DLog 证明。

本方作为 Receiver，跟其他参与方 $j$ 调用子协议 Endemic OT。
输入 $\boldsymbol{\beta}_{i,j} \in \left\{0,1\right\}^\kappa$，即 Base OT 选择位向量。
输出 $R_{i,j}^0, R_{i,j}^1$，是子协议中途产生的通信消息。

【交换】
\begin{itemize}
\item
  Shamir 相关的消息，即 $\mathbf{F}_i, \varepsilon_{1,i}$ 以及 DLog 证明。
\item
  Endemic OT Receiver 消息，即 $R_{i,j}^0, R_{i,j}^1$。
\item
  揭示“零和掩码分片的承诺”：$\epsilon^\circ_{i,j},\nu_{i,j}$。
\end{itemize}

【检查-K】本方按 \eqref{eq:orch:feld-com} 重新计算 $\mathrm{Com}_{1,j}$，
检查与 \hyperref[sec:orchestration:round1]{第一轮} 所收到的是否相等。
本方按公式 \eqref{eq:orch:seed-com}，交换 $i,j$ 后重算 $\mathrm{Com}_{\epsilon,j,i}$，
检查对方揭示的 $\epsilon^\circ_{j,i},\nu_{j,i}$
与 \hyperref[sec:orchestration:round1]{第一轮} 所收到的承诺一致。
以及，检查所有 DLog 证明。

【检查-K】所有检查通过后，本方将两个贡献逐位异或，得到公式 \eqref{eq:orch:symmetric-blind} 的共享种子。
双方按同一公式计算，因而各自保存的 $\epsilon_{i,j}$ 与 $\epsilon_{j,i}$ 相等。
\begin{equation}
  \label{eq:orch:symmetric-blind}\tag{sym-blind}
  \epsilon_{i,j}:=\epsilon^\circ_{i,j}\oplus\epsilon^\circ_{j,i}
  =\epsilon_{j,i}.
\end{equation}

先承诺再揭示，使一方无法在看到对方贡献后再改变自己的贡献，以指定合成结果。

【检查-K】恶意方仍可能拒绝揭示；拒绝揭示、揭示长度不符或承诺验证失败时，执行 K 类处理。
共享种子随成功生成的 Keyshare 持久保存，后续 PreSign 使用它派生各次会话的掩码。
这一初始化方式对应 DKLS23 第 3.1 节中 $F_{\mathrm{Zero}}$ 的构造说明。

最后，把各方揭示的 Feldman 向量加起来，得到全局多项式承诺 $\mathbf{F}$：
\begin{equation}
\label{eq:orch:feld-agg}\tag{F-agg}
\mathbf{F} := \sum_{j\in\mathbb{K}} \mathbf{F}_j
= \Bigl(\textstyle\sum_j c_{j,0}\cdot G,\ \cdots,\ \sum_j c_{j,t-1}\cdot G\Bigr).
\end{equation}

这就是全局多项式 $P:=\sum_j P_j$ 的群承诺。
记“在指数上求值”为 $\mathbf{F}(z) := \sum_{w\in[0,t)} z^w \cdot \mathbf{F}[w]$，
则 $\mathbf{F}(z) = P(z)\cdot G$。私钥为常数项 $P(0)$，公钥为常数项 $\mathbf{F}(0)$。

\subsubsection{Endemic OT 收尾，PPRF 建树，Shamir 划分私钥}

本方计算每个参与方对应的 Shamir 私钥份额 $d_{i,j}$，
如公式 \eqref{eq:orch:shamir-share}，其中 $d_{i,i}$ 在本地保留。
\begin{equation}
\label{eq:orch:shamir-share}\tag{shamir}
d_{i,j} := P_i(j), ~\forall j\in\mathbb{K}.
\end{equation}

本方作为 Sender，跟其他方调用子协议 Endemic OT，只执行到 Sender 备好消息。
输入为 $R_{j,i}^0, R_{j,i}^1$。
输出有
(i) Base OT 密钥对 $\rho^0_{\ell,j,i}, \rho^1_{\ell,j,i}, \ell\in[0, \kappa)$；
(ii) Endemic OT Sender 消息 $M_{\ell,j,i}^0, M_{\ell,j,i}^1$。

本方作为 Sender，跟其他方执行子协议 PPRF，只执行 Build 和 Prove。
本方跟每一方 $j$ 得到 $\kappa/k$ 棵 PPRF 树，
第 $h$ 棵树 $\mathcal{T}_{h,i,j}$ 有：
\begin{itemize}
\item
  叶子节点的集合 $\mathcal{A}_{h,i,j}$ 如公式 \eqref{eq:orch:all-leaves}。
  \footnote{记号 $\mathcal{A}$ 取自法语 arbre。
  这样取名是因为，L（对应英语 leaf）和 F（对应法语 feuille）都已经被高频地占用。}
  \begin{equation}
    \tag{all-leaves}\label{eq:orch:all-leaves}
    \mathcal{A}_{h,i,j} := \left\{\left(\mathcal{T}_{h,i,j}\right)^k_z ~:~ z\in[0, 2^k)\right\}.
  \end{equation}
\item
  除第一层以外，每一层有选 0 修正值和选 1 修正值。
  所有选 0 修正值按层拼接为向量 $t_{h,i,j}^0$，
  所有选 1 修正值按层拼接为向量 $t_{h,i,j}^1$。
\item
  证明材料 $\tilde t_{h,i,j}, \tilde s_{h,i,j}$。
\end{itemize}

本方对其他参与方、对树编号 $h \in [0, \kappa/k)$
以及对 Base OT 实例编号 $\ell \in [0, \kappa)$，交换：
\begin{itemize}
\item
  Shamir 相关的消息：$d_{i,j}, \mathbf{F}$。
\item
  Endemic OT Sender 消息：$M_{\ell,j,i}^0, M_{\ell,j,i}^1$。
\item
  PPRF Sender 消息：$t_{h,i,j}^0, t_{h,i,j}^1, \tilde t_{h,i,j}, \tilde s_{h,i,j}$。
\end{itemize}

【检查-K】对其他参与方进行如公式
\eqref{eq:orch:echo-check} 和 \eqref{eq:orch:share-check} 的验证。
公式 \eqref{eq:orch:echo-check} 的用意是确保各方产生一致的 $\mathbf{F}$，
从而产生一致的公钥；式中 $\mathbf{F}^{(j)}$ 是参与方 $j$ 发来的聚合承诺；
$\mathbf{F}$ 是按公式 \eqref{eq:orch:feld-agg} 本地算出的聚合承诺。
公式 \eqref{eq:orch:share-check} 的用意是检查 $j$ 发来的 Shamir 份额的诚实性，
确保份额由他所承诺的多项式产生。
\begin{align}
\label{eq:orch:echo-check}\tag{echo}
\mathbf{F}^{(j)} &\stackrel{?}{=}
\mathbf{F} = \sum_{w\in\mathbb{K}} \mathbf{F}_w, \\
\label{eq:orch:share-check}\tag{share}
d_{j,i} \cdot G &\stackrel{?}{=}
\mathbf{F}_j(i) = \sum_{w\in[0,t)} i^w \cdot \mathbf{F}_j[w].
\end{align}

把两个子协议 Endemic OT 和 PPRF 执行完。
本方首先作为 Receiver 与其他各方继续执行 Endemic OT。
对每个 $\ell\in[0, \kappa)$，
输入 $R_{i,j}^0, R_{i,j}^1, M_{\ell,i,j}^0, M_{\ell,i,j}^1$，
输出 Base OT 密钥 $\rho^{\beta_{\ell,i,j}}_{\ell,i,j}$。

本方然后作为 Receiver 与其他各方继续执行 PPRF。
对每棵 PPRF 树编号 $h\in[0, \kappa/k)$，
输入对应的 $k$ 个 Base OT 选择位与单侧密钥
$\bigl(\beta_{\ell,i,j},\rho^{\beta_{\ell,i,j}}_{\ell,i,j}\bigr)_{\ell\in[hk,(h+1)k)}$，
以及收到的修正向量 $t^0_{h,j,i},t^1_{h,j,i}$ 和证明材料 $\tilde t_{h,j,i},\tilde s_{h,j,i}$。
输出打孔位置 $y_{h,i,j}$ 与打孔叶子 $\mathcal{A}^*_{h,i,j}$。

【检查-K】如果公式 \eqref{eq:orch:echo-check}、\eqref{eq:orch:share-check}，
或子协议 Endemic OT、PPRF 中的检验失败，则执行 K 类处理。
如果这些检验都通过，则聚合生成本方 Shamir 份额 $x_i$ 如公式 \eqref{eq:orch:xi}。
\begin{equation}
  \label{eq:orch:xi}
  \tag{xi}
  x_i := \sum_{j\in\mathbb{K}} d_{j,i} = P(i).
\end{equation}

\subsubsection{公开最终份额}

本方获取公钥 $Y := \mathbf{F}(0)$，即向量 $\mathbf{F}$ 的首项。
计算 $S_i := x_i \cdot G$ 并出具相应的 DLog 证明 $\pi_i$。

交换 $S_i$，$\pi_i$。

【检查-K】然后检验所有 DLog 证明 $\pi_j$，以及用拉格朗日插值重构公钥，
如公式 \eqref{eq:orch:lagrange-interp} 等号右边，本次使用全体 Keygen 参与方 $\mathbb K$。
后文对签名集合 $\mathbb S$ 插值时，将系数定义中的 $\mathbb K$ 替换为 $\mathbb S$。
\begin{align}
  \label{eq:orch:lagrange-coef}\tag{coef}
  \lambda(j,\mathbb{K}) &:= \prod_{k\in\mathbb{K}\setminus\{j\}} \frac{-k}{j - k}, \\
  \label{eq:orch:lagrange-interp}\tag{interp}
  Y &\phantom{:}\stackrel{?}{=} \sum_{j\in\mathbb{K}}\lambda(j,\mathbb{K}) \cdot S_j.
\end{align}

\begin{center}\rule{0.5\linewidth}{0.5pt}\end{center}

\section{PreSign}

目标：把跟消息 $m$ 无关的所有协议步骤一次性跑完，
产出每方的 PreSignature。PreSign 交换两轮消息，之后绑定消息 $m$ 的 Sign 再交换一轮。
初始化完成后，整个签名流程共三轮；第一轮同时发送 nonce 承诺与 SoftSpoken Receiver 消息。

【检查-S】ECDSA 签名公式回顾（伪代码，全局视角）：

\begin{codebox}
\Procname{$\proc{Sign}(m)$}
\li $k \gets \sum_{i \in\mathbb{S}} r_i$
\li $\phi \gets \sum_{i \in\mathbb{S}} \phi_i$
\li $R \gets k \cdot G$
\li 若 $R=\mathcal O$，则 \Return $\bot$
\li $r \gets R.\mathrm{x} \bmod n$
\li 若 $r=0$ 或 $k\phi=0 \pmod n$，则 \Return $\bot$
\li $s \gets (k\phi)^{-1}\bigl(m\phi + r\,x'\phi\bigr)$
\Comment 乘积 $k\phi$、$x'\phi$ 由 MtA 兑现为加法份额
\li 若 $s=0$，则 \Return $\bot$
\li \Return $(r, s)$
\end{codebox}

【检查-S】其中 $x' := x + \nabla x$，标量运算均在 $\mathbb Z_n$ 中进行。
$\mathcal O$ 表示椭圆曲线的无穷远点，$\bot$ 表示本次尝试失败。
返回失败时须废弃本次预签名状态；具体检查位置及重试规则见下文。

\subsection{规格}

【输入】参与方 $i$ 输入 Keyshare、签名者集合 $\mathbb{S}$，BIP-32 私钥偏移量 $\nabla x$。
无需消息哈希 $m$，其在 Sign 阶段绑定。

【输出】每方 $i$ 获得预签名结构体，如公式 \eqref{eq:orch:presig}。
其中 $r$ 是 ECDSA 签名的 $r$ 字段，全员一致；
$s_{1,i}$ 已经是最终值，待 Sign 阶段直接广播；
$s_{0,i}^\circ$ 是未绑 $m$ 的中间值，待 Sign 阶段叠加 $m\cdot\phi_i$；
$\phi_i$ 本方私有，用于 Sign 时绑定一个 $m$。
\begin{equation}
\label{eq:orch:presig}\tag{presig}
\mathrm{PreSignature}_i := \bigl(r,\, s_{1,i},\, s_{0,i}^\circ,\, \phi_i\bigr).
\end{equation}

【安全承诺】

\begin{itemize}
\item
  秘密分片 $r_i$，$x'_i$，$\phi_i$ 不泄露。
\item
  MtA 份额在不知情者看来是均匀随机的。
\item
  【检查-P】$\Gamma$ 一致性检查把各方喂给 RVOLE 的输入与广播的 $R_j, Y_j$ 绑住，
  transcript hash $d$ 用于核对全员视图；失败时执行 P 类处理。
\end{itemize}

【安全警告】PreSignature 应当用后即焚。
若同一份 PreSignature 在同一派生私钥 $x'$ 下，对 $m_1\ne m_2\pmod n$ 产生两份有效签名，
则两份签名复用了同一个 nonce $k$，并有相同的非零 $r$。
对未经 $s\mapsto n-s$ 规范化的签名值 $s^{(1)},s^{(2)}$，任意获得两份签名的人都能恢复
\begin{equation}
  k=(m_1-m_2)(s^{(1)}-s^{(2)})^{-1},\qquad
  x'=(s^{(1)}k-m_1)r^{-1}\pmod n.
\end{equation}
这是由 $s^{(a)}k=m_a+rx'$ 直接相减得到的；上述条件保证两个分母均非零。
若签名经过 low-S 规范化，可枚举两个 $s$ 的正负号，并用公钥核对候选私钥。
仅从重复发送的签名份额恢复 $\phi_i$，不足以直接断言已恢复本方 $x'_i$；这里采用完整私钥恢复的结论。

\subsection{实施}

\subsubsection{得到 MtA 乘法 $r$-分片、$x'$-分片与加法 $y$-分片}

获取随机 nonce 分片 $r_i \stackrel{\$}{\leftarrow} \mathbb{Z}_n^*$、
MtA 盲化分片 $\phi_i \stackrel{\$}{\leftarrow} \mathbb{Z}_n^*$。

计算衍生分片 $x'_i$，如公式 \eqref{eq:orch:xi-derived}。
式中 $\lambda(i,\mathbb{S})\cdot x_i$ 是拉格朗日插值后的私钥分片；
$\zeta_i$ 是由 Keyshare 中的共享种子派生的零和掩码，
详见公式 \eqref{eq:orch:zeta-pair} 和 \eqref{eq:orch:zeta}；
$\nabla x$ 是 BIP-32 偏移量，系数 $|\mathbb{S}|^{-1}$ 把偏移量均摊到每个签名方。
\begin{equation}
  \label{eq:orch:xi-derived}\tag{xi-derived}
  x'_i := \lambda(i,\mathbb{S})\cdot x_i + \zeta_i + \nabla x \cdot |\mathbb{S}|^{-1}.
\end{equation}

协议 Keygen 已按公式 \eqref{eq:orch:symmetric-blind} 为每对参与方。
本次 PreSign 中，对每个 $j\in\mathbb{S}\setminus\{i\}$，
本方用 $\epsilon_{i,j}$ 与公开的会话信息派生同一个标量 $\zeta_{i,j}=\zeta_{j,i}$，
如公式 \eqref{eq:orch:zeta-pair}。

\begin{equation}
  \label{eq:orch:zeta-pair}\tag{zeta-pair}
  \begin{aligned}
    I, J &:= \min(i,j),\max(i,j); \\
    \zeta_{i,j}&:=\mathrm{Hash}\bigl(
      \texttt{"zero\_mask"},\mathrm{sid},\mathrm{sort}(\mathbb{S}),
        I, J,\epsilon_{I, J}\bigr) \\
    &\in\mathbb{Z}_n.
  \end{aligned}
\end{equation}

式中 $\mathtt{sort}$ 是升序排序函数，确保 $\zeta_{i,j}$ 与 $\mathbb{S}$ 的顺序无关。
该式引入参数 $\mathrm{sid}$，意味着 $\zeta_{i,j}$ 不可跨会话复用。

本方按编号大小对这些标量取正负号，求和得到 $\zeta_i$：
\begin{equation}
  \label{eq:orch:zeta}\tag{zeta}
  \zeta_i := \sum_{\substack{j\in\mathbb{S}\\j<i}} \zeta_{i,j}
  - \sum_{\substack{j\in\mathbb{S}\\j>i}} \zeta_{i,j}.
\end{equation}
对每对参与方，编号较大的一方加上 $\zeta_{i,j}$，编号较小的一方减去同一个值。
因此，全员求和时各项两两抵消，保证 $\sum_{i\in\mathbb{S}}\zeta_i=0$，
从而再随机化不改变各方衍生私钥分片的总和。
派生过程只依赖已有的共享种子和双方一致的公开输入，所以 PreSign 无需为零和掩码增加通信。

对 $R_i := r_i \cdot G$ 做哈希承诺如公式 \eqref{eq:orch:com-Ri}，
式中 $\varepsilon_{R,i} \stackrel{\$}{\leftarrow} \left\{0,1\right\}^{256}$ 为盲化项。
\begin{equation}
  \label{eq:orch:com-Ri}\tag{Com-Ri}
  \mathrm{Com}_{R,i} := \mathrm{Hash}(\mathrm{sid}, i, R_i, \varepsilon_{R,i}).
\end{equation}

本方作为 Receiver，与其他方调用子协议 SoftSpoken。
输入 Keyshare 中的完整叶子 $\mathcal{A}_{h,i,j}$，其中 $h\in[0,\kappa/k)$，
以及随机选择串 $\boldsymbol{\beta}_{j,i} \stackrel{\$}{\leftarrow} \left\{0,1\right\}^L$。
生成待发送的消息 $u_{j,i}, \tilde{\beta}_{j,i}, \tau_{j,i}$。
本地输出并保存单侧 OT 密钥 $\rho^{\beta_{\ell,j,i}}_{\ell,j,i}$，其中 $\ell\in[0,L)$，
供后续 RVOLE Receiver 使用。

计算向量 $\boldsymbol{\beta}_{j,i}$ 的 gadget 加权和，
得到 $\beta_{j,i}:=\langle \boldsymbol{g}, \boldsymbol{\beta}_{j,i}\rangle \in \mathbb{Z}_n$。

【交换】第一轮，同时发送 $\mathrm{Com}_{R,i}$ 与 $u_{j,i}, \tilde{\beta}_{j,i}, \tau_{j,i}$。
SoftSpoken Receiver 消息的生成不依赖对端承诺，因此这两类消息可以同轮发送。
收齐第一轮消息后继续；承诺的揭示留到第二轮。

本地计算全体承诺 $\mathrm{Com}_{R,j}, j\in\mathbb{S}$ 的哈希 $d$ 如公式 \eqref{eq:orch:com-hash}，
其将在第二轮维护各方的一致性。式中 $Y' := Y+\nabla x \cdot G$ 为衍生公钥。
\begin{equation}
  \label{eq:orch:com-hash}\tag{com-hash}
  d := \mathrm{Hash}\left(
    \mathrm{sid}, Y', \left\{\mathrm{Com}_{R,j} : j \in \mathbb{S}\right\}
  \right).
\end{equation}

【检查-P】然后，本方作为 Sender，与其他方调用子协议 SoftSpoken。
输入 Keyshare 中的打孔位置与打孔叶子 $\bigl(y_{h,i,j}, \mathcal{A}^*_{h,i,j}\bigr)$，
以及收到的 $u_{i,j}, \tilde{\beta}_{i,j}, \tau_{i,j}$。
验证通过后，输出 $L$ 对 OT 密钥 $\rho^0_{\ell,i,j}, \rho^1_{\ell,i,j}$，
其中 $h\in[0,\kappa/k)$，$\ell\in[0,L)$。验证失败时，对该对端 $j$ 执行 P 类处理。

【检查-P】然后，本方作为 Sender，与其他方调用子协议 RVOLE，详见「随机 VOLE」。
秘密输入为 $(r_i,x'_i)$，并使用上述 OT 密钥对。
输出修正矩阵 $\tilde a_{i,j}$、响应 $\eta_{i,j}$ 、哈希校验 $\mu_{i,j}$；
以及 MtA 加法分片 $y^\mathtt{r}_{i,j}, y^\mathtt{x}_{i,j}$。
这里生成校验材料，由对端的 Receiver 执行校验；失败处理见下文 RVOLE Receiver 步骤。
\begin{equation}
  \label{eq:orch:rand-mta}\tag{rand-mta}
  \begin{aligned}
    y^\mathtt{r}_{i,j} + z^\mathtt{r}_{i,j} &:= r_i \cdot \beta_{i,j}, \\
    y^\mathtt{x}_{i,j} + z^\mathtt{x}_{i,j} &:= x'_i \cdot \beta_{i,j}.
  \end{aligned}
\end{equation}
两个加法分片满足公式 \eqref{eq:orch:rand-mta}。
其中 $\beta_{i,j}$ 是 RVOLE Receiver 也就是参与方 $j$ 的私有乘法分片；
$z$-分片是他的私有加法分片，此时还没出现。

\subsubsection{得到 MtA 加法 $z$-分片}

算 $\Gamma$ 一致性点 $\Gamma^\mathtt{r}_{i,j} := y^\mathtt{r}_{i,j} \cdot G$，
$\Gamma^\mathtt{x}_{i,j} := y^\mathtt{x}_{i,j} \cdot G$。
这两项用于把 RVOLE 的输入绑到将要交换的 $R_i, Y_i$ 上。

算 $\phi$ 专用的偏移量 $\psi_{i,j} := \phi_i - \beta_{j,i}$。

【交换】第二轮，发送以下消息。
\begin{itemize}
\item
  RVOLE Sender 消息 $\tilde a_{i,j}, \eta_{i,j}, \mu_{i,j}$；
\item
  承诺 $\mathrm{Com}_{R,i}$ 的揭示，即 $R_i, \varepsilon_{R,i}$；
\item
  本方公钥分片 $Y_i$；
\item
  一致性约束 $\Gamma^\mathtt{r}_{i,j}, \Gamma^\mathtt{x}_{i,j}$；
\item
  偏移量 $\psi_{i,j}$；
\item
  一致性约束 $d$，详见公式 \eqref{eq:orch:com-hash}。
\end{itemize}

【检查-P】重新计算 $\mathrm{Com}_{R,j}$，检查与第一轮所收到的是否相等。
揭示失败时，对该对端 $j$ 执行 P 类处理。

【检查-P】检查各方发来的 $d$ 与本地计算的 $d$ 是否相等。
若不一致，按无法定位的情形执行 P 类处理，取 $B=\mathbb S\setminus\{i\}$。

【检查-P】本方作为 Receiver，与其他方调用子协议 RVOLE。
输入本方的选择串 $\boldsymbol{\beta}_{j,i}$、单侧 OT 密钥
$\bigl(\rho^{\beta_{\ell,j,i}}_{\ell,j,i}\bigr)_{\ell\in[0,L)}$，
以及 $j$ 发来的 $\tilde a_{j,i}, \eta_{j,i}, \mu_{j,i}$，
子协议须先通过 $\mu_{j,i}$ 的校验，才输出 MtA 加法分片 $z^\mathtt{r}_{j,i}, z^\mathtt{x}_{j,i}$。
若子协议报告失败，则对该对端 $j$ 执行 P 类处理。

【检查-P】检查 $\Gamma$ 一致性如下式；任一等式失败时，对相应对端 $j$ 执行 P 类处理。
RVOLE 协议本身不保证输入的乘法分片是下游协议所需的。
这就需要靠 $\Gamma$ 把 RVOLE 的输入 $r_j, x'_j$ 跟此前广播的 $R_j, Y_j$ 关联起来。
\begin{align*}
\beta_{j,i} \cdot R_j &\stackrel{?}{=} \Gamma^\mathtt{r}_{j,i} + z^\mathtt{r}_{j,i} \cdot G, \\
\beta_{j,i} \cdot Y_j &\stackrel{?}{=} \Gamma^\mathtt{x}_{j,i} + z^\mathtt{x}_{j,i} \cdot G.
\end{align*}

【检查-P】检查派生公钥是否一致，即 $\sum_{j\in\mathbb{S}}Y_j \stackrel{?}{=} Y'$，
其中 $Y_j:=x'_j G$。若不一致，取 $B=\mathbb S\setminus\{i\}$，执行 P 类处理。

【检查-S】先计算总点 $R := \sum_{j\in\mathbb S} R_j$，检查 $R\ne\mathcal O$。
若 $R=\mathcal O$，则中止本次 PreSign，废弃本次临时状态，不能取其横坐标。
检查通过后计算 $r := R.\mathrm{x} \bmod n$，取值范围为 $0\le r<n$。
若 $r=0$，同样中止并废弃本次状态；只有 $1\le r<n$ 时才继续生成 PreSignature。
各方独立采样的 $r_i$ 均非零，并不能保证总 nonce $k=\sum_{j\in\mathbb S}r_j$ 非零。

然后计算 PreSignature 的两个分量，如公式 \eqref{eq:orch:presig-steps}。
\begin{equation}
\label{eq:orch:presig-steps}\tag{presig-steps}
\begin{aligned}
  \Phi_i &:= \phi_i + \sum_{j\ne i}\psi_{j,i}, \\
  V^\mathtt{r}_i &:= \sum_{j\ne i}(y^\mathtt{r}_{i,j} + z^\mathtt{r}_{j,i}), \\
  V^\mathtt{x}_i &:= \sum_{j\ne i}(y^\mathtt{x}_{i,j} + z^\mathtt{x}_{j,i}), \\
  s_{1,i} &:= r_i \cdot \Phi_i + V^\mathtt{r}_i, \\
  s_{0,i}^\circ &:= r \cdot (x'_i \cdot \Phi_i + V^\mathtt{x}_i).
\end{aligned}
\end{equation}

至此凑齐 PreSignature 的全部分量，见公式 \eqref{eq:orch:presig}。

\begin{center}\rule{0.5\linewidth}{0.5pt}\end{center}

\section{Sign 1 轮}

目标：用消息哈希 $m$ 绑定 PreSignature，输出 ECDSA 签名 $(r, s)$。

\subsection{规格}

【输入】
\begin{itemize}
\item
  各方的 PreSignature；
\item
  各方相同的消息哈希 $m \in \mathbb{Z}_n$。
\end{itemize}

【输出】标准 ECDSA 签名 $(r, s)$，全员一致。

【安全承诺】广播的 $(s_{0,i}, s_{1,i})$ 不泄露 $\phi_i$ 与私钥分片，
前提是「PreSignature 用完即弃」。

\subsection{实施}

【检查-S】本方先检查输入 PreSignature 中的 $r$ 满足 $1\le r<n$，否则废弃该预签名并中止。
本方计算 $s_{0,i} := s_{0,i}^\circ + m \cdot \phi_i$。
在发送签名份额之前，将该 PreSignature 标记为已使用，然后在第三轮交换 $s_{0,i}, s_{1,i}$。
无论本次 Sign 成功还是失败，该预签名均不得再次使用。

【检查-S】收齐份额后，在 $\mathbb Z_n$ 中计算分母 $D := \sum_{j\in\mathbb S}s_{1,j}$。
若 $D=0$，则中止本次 Sign，废弃该预签名，不能继续求逆。
正常执行时 $D=k\phi$，其中 $\phi=\sum_{j\in\mathbb S}\phi_j$；
即使 PreSign 已排除 $k=0$，各方非零的 $\phi_j$ 仍可能相加为零。
此处检查的是份额的总和，单个 $s_{1,j}$ 或 $s_{0,j}$ 为零不构成失败条件。
只有 $D\ne0$ 时，才计算 ECDSA $s$-字段，如公式 \eqref{eq:orch:ecdsa-s}。
\begin{equation}
  \label{eq:orch:ecdsa-s}\tag{ecdsa-s}
  s := D^{-1}\sum_{j\in\mathbb S}s_{0,j} \pmod n.
\end{equation}

【检查-S】将 $s$ 表示为 $0\le s<n$ 的整数；若 $s=0$，则中止并废弃该预签名。
只有 $1\le r,s<n$ 时，才以消息哈希 $m$ 和派生公钥 $Y'$ 进行 ECDSA 标准验签。
验签失败则中止并废弃该预签名；验签通过才输出 $(r,s)$，随后销毁本次预签名的秘密状态。

【检查-S】上述 Sign 检查失败时，本次尝试均不输出签名。
若需重试，先确认签名集合未涉及已停用对端，再使用新的 PreSign 会话号。
重新采样 $r_i,\phi_i$，并重新执行完整的 PreSign。
本次产生的 OT 扩展材料、MtA 份额及 PreSignature 均随本次状态一起废弃，不能带入重试。
